\documentclass[10pt]{article}
\usepackage[utf8]{inputenc}
\usepackage[T2A]{fontenc}
\usepackage[english, russian]{babel}
\usepackage{bbold}
\usepackage{amsfonts}
\usepackage{amsmath}
\usepackage{amssymb}
\usepackage{amsthm}
\usepackage{mathtools}


\newtheorem{theorem}{Теорема}
\newtheorem{lemma}{Лемма}
\newtheorem{definition}{Определение}
\newtheorem*{stat}{Утверждение}
\newtheorem*{remark}{Замечание}
\newtheorem*{property}{Свойство}
\newtheorem{corollary}{Следствие}[theorem]
\newtheorem{example}{Пример}

\newcommand{\Bb}{\mathcal{B}}
\newcommand{\A}{\mathcal{A}}
\newcommand{\B}{\mathcal{B}}
\newcommand{\E}{\mathcal{E}}
\newcommand{\0}{\mathbf{0}}
\newcommand{\R}{\mathbb{R}}
\newcommand{\N}{\mathbb{N}}
\newcommand{\Z}{\mathbb{Z}}
\newcommand{\Q}{\mathbb{Q}}
\newcommand{\CC}{\mathbb{C}}
\newcommand{\Ker}{\operatorname{Ker}}
\newcommand{\rg}{\operatorname{rg}}
\newcommand{\tr}{\operatorname{tr}}
\newcommand{\diag}{\operatorname{diag}}
\newcommand{\rank}{\operatorname{rank}}

\begin{document}

\begin{definition}
\textbf{Полем} называется множество \(\mathbb{F}\), на котором определены две бинарные операции: сложение (\(+\)) и умножение (\(\cdot\)), удовлетворяющие следующим свойствам (аксиомам поля):
\begin{enumerate}
    \item \textbf{Ассоциативность сложения:} \(\forall a,b,c \in \mathbb{F}: (a+b)+c = a+(b+c)\).
    \item \textbf{Коммутативность сложения:} \(\forall a,b \in \mathbb{F}: a+b = b+a\).
    \item \textbf{Существование нейтрального элемента по сложению (нуля):} \(\exists 0 \in \mathbb{F} : \forall a \in \mathbb{F}, a+0 = a\).
    \item \textbf{Существование противоположного элемента:} \(\forall a \in \mathbb{F}, \exists (-a) \in \mathbb{F} : a+(-a) = 0\).
    \item \textbf{Ассоциативность умножения:} \(\forall a,b,c \in \mathbb{F}: (a\cdot b)\cdot c = a\cdot(b\cdot c)\).
    \item \textbf{Коммутативность умножения:} \(\forall a,b \in \mathbb{F}: a\cdot b = b\cdot a\).
    \item \textbf{Существование нейтрального элемента по умножению (единицы):} \(\exists 1 \in \mathbb{F}, 1 \neq 0 : \forall a \in \mathbb{F}, a\cdot 1 = a\).
    \item \textbf{Существование обратного элемента:} \(\forall a \neq 0 \in \mathbb{F}, \exists a^{-1} \in \mathbb{F} : a \cdot a^{-1} = 1\).
    \item \textbf{Дистрибутивность:} \(\forall a,b,c \in \mathbb{F}: a\cdot(b+c) = a\cdot b + a\cdot c\).
\end{enumerate}
\end{definition}

Наиболее часто используются поля действительных чисел \(\mathbb{R}\) и комплексных чисел \(\mathbb{C}\). Элементы поля называются \textit{скалярами}.

\section{Линейное (векторное) пространство}

\begin{definition}
\textbf{Линейным (или векторным) пространством} над полем \(\mathbb{F}\) называется множество \(V\), на котором определены две операции:
\begin{enumerate}
    \item[\textbf{1.}] \textit{Внутренняя операция:} сложение векторов, ставящее в соответствие любой упорядоченной паре элементов \(\mathbf{x}, \mathbf{y} \in V\) элемент \(\mathbf{x} + \mathbf{y} \in V\).
    \item[\textbf{2.}] \textit{Внешняя операция:} умножение вектора на скаляр, ставящее в соответствие любому скаляру \(\lambda \in \mathbb{F}\) и любому вектору \(\mathbf{v} \in V\) элемент \(\lambda \mathbf{v} \in V\).
\end{enumerate}
Эти операции должны удовлетворять следующим аксиомам для любых \(\mathbf{u}, \mathbf{v}, \mathbf{w} \in V\) и любых \(\alpha, \beta \in \mathbb{F}\):

\begin{enumerate}
    \item \(\mathbf{u} + (\mathbf{v} + \mathbf{w}) = (\mathbf{u} + \mathbf{v}) + \mathbf{w}\) (ассоциативность сложения).
    \item \(\mathbf{u} + \mathbf{v} = \mathbf{v} + \mathbf{u}\) (коммутативность сложения).
    \item Существует элемент \(\mathbf{0} \in V\) (называемый \textit{нулевым вектором}) такой, что \(\mathbf{v} + \mathbf{0} = \mathbf{v}\) для любого \(\mathbf{v} \in V\).
    \item Для каждого вектора \(\mathbf{v} \in V\) существует вектор \(-\mathbf{v} \in V\) (называемый \textit{противоположным}) такой, что \(\mathbf{v} + (-\mathbf{v}) = \mathbf{0}\).
    \item \(\alpha(\beta\mathbf{v}) = (\alpha\beta)\mathbf{v}\) (ассоциативность умножения на скаляр).
    \item \(1 \cdot \mathbf{v} = \mathbf{v}\) (унитарность).
    \item \((\alpha + \beta)\mathbf{v} = \alpha\mathbf{v} + \beta\mathbf{v}\) (дистрибутивность относительно сложения скаляров).
    \item \(\alpha(\mathbf{u} + \mathbf{v}) = \alpha\mathbf{u} + \alpha\mathbf{v}\) (дистрибутивность относительно сложения векторов).
\end{enumerate}
Элементы множества \(V\) называются \textit{векторами}.
\end{definition}

\begin{example}
Примерами линейных пространств являются:
\begin{itemize}
    \item Множество \(\mathbb{R}^n\) (векторы-строки или столбцы) над полем \(\mathbb{R}\).
    \item Множество всех матриц размера \(m \times n\) с элементами из \(\mathbb{F}\) (обозначается \(\mathbb{F}^{m \times n}\)).
    \item Множество всех многочленов от одной переменной с коэффициентами из \(\mathbb{F}\) (обозначается \(\mathbb{F}[x]\)).
    \item Множество всех функций, определенных на некотором множестве \(X\) и принимающих значения в \(\mathbb{F}\).
\end{itemize}
\end{example}

\subsection{Простейшие свойства линейного пространства}

Из аксиом линейного пространства можно вывести ряд фундаментальных свойств.

\begin{theorem}
Для любого вектора \(\mathbf{v} \in V\) и любого скаляра \(\lambda \in \mathbb{F}\) выполняются следующие утверждения:
\begin{enumerate}
    \item \(0 \cdot \mathbf{v} = \mathbf{0}\).
    \item \(\lambda \cdot \mathbf{0} = \mathbf{0}\).
    \item Если \(\lambda\mathbf{v} = \mathbf{0}\), то либо \(\lambda = 0\), либо \(\mathbf{v} = \mathbf{0}\).
    \item \((-1) \cdot \mathbf{v} = -\mathbf{v}\).
\end{enumerate}
\end{theorem}

\begin{proof}
Докажем каждое свойство.
\begin{enumerate}
    \item Рассмотрим вектор \(\mathbf{v}\). Используя аксиомы, получаем:
    \[
    0\cdot\mathbf{v} = (0+0)\cdot\mathbf{v} = 0\cdot\mathbf{v} + 0\cdot\mathbf{v}.
    \]
    Прибавим к обеим частям этого равенства вектор \(-(0\cdot\mathbf{v})\), противоположный к \(0\cdot\mathbf{v}\):
    \[
    0\cdot\mathbf{v} + (-(0\cdot\mathbf{v})) = (0\cdot\mathbf{v} + 0\cdot\mathbf{v}) + (-(0\cdot\mathbf{v})).
    \]
    Левая часть равна \(\mathbf{0}\) по аксиоме 4. В правой части, используя ассоциативность (аксиома 1), получаем \(0\cdot\mathbf{v} + (0\cdot\mathbf{v} + (-(0\cdot\mathbf{v}))) = 0\cdot\mathbf{v} + \mathbf{0} = 0\cdot\mathbf{v}\). Следовательно, \(\mathbf{0} = 0\cdot\mathbf{v}\).

    \item Аналогично, \(\lambda \cdot \mathbf{0} = \lambda \cdot (\mathbf{0} + \mathbf{0}) = \lambda\mathbf{0} + \lambda\mathbf{0}\). Прибавляя к обеим частям \(-\lambda\mathbf{0}\), получаем \(\mathbf{0} = \lambda\mathbf{0}\).

    \item Пусть \(\lambda\mathbf{v} = \mathbf{0}\) и \(\lambda \neq 0\). Так как \(\lambda\) -- ненулевой элемент поля, для него существует обратный \(\lambda^{-1}\). Умножим обе части равенства на \(\lambda^{-1}\):
    \[
    \lambda^{-1}(\lambda\mathbf{v}) = \lambda^{-1}\mathbf{0}.
    \]
    Левая часть: \((\lambda^{-1}\lambda)\mathbf{v} = 1\cdot\mathbf{v} = \mathbf{v}\) (по аксиомам 5 и 6). Правая часть: \(\lambda^{-1}\mathbf{0} = \mathbf{0}\) (по свойству 2, уже доказанному). Следовательно, \(\mathbf{v} = \mathbf{0}\).

    \item Рассмотрим сумму \((-1)\cdot\mathbf{v} + \mathbf{v}\):
    \[
    (-1)\cdot\mathbf{v} + \mathbf{v} = (-1)\cdot\mathbf{v} + 1\cdot\mathbf{v} = (-1 + 1)\cdot\mathbf{v} = 0\cdot\mathbf{v} = \mathbf{0}.
    \]
    Так как вектор \((-1)\cdot\mathbf{v}\) в сумме с \(\mathbf{v}\) дает \(\mathbf{0}\), по определению он является противоположным для \(\mathbf{v}\), т.е. \((-1)\cdot\mathbf{v} = -\mathbf{v}\).
\end{enumerate}
\end{proof}

\section{Линейное подпространство}

\begin{definition}
Непустое подмножество \(L \subseteq V\) линейного пространства \(V\) над полем \(\mathbb{F}\) называется \textbf{линейным подпространством} (или просто подпространством) пространства \(V\), если оно само является линейным пространством относительно операций сложения и умножения на скаляр, определенных в \(V\).
\end{definition}

\subsection{Критерий подпространства}

\begin{theorem}[Критерий подпространства]
Непустое подмножество \(L \subseteq V\) линейного пространства \(V\) над полем \(\mathbb{F}\) является подпространством тогда и только тогда, когда оно замкнуто относительно линейных операций, то есть выполняются два условия:
\begin{enumerate}
    \item \(\forall \mathbf{x}, \mathbf{y} \in L: \mathbf{x} + \mathbf{y} \in L\) (замкнутость относительно сложения).
    \item \(\forall \lambda \in \mathbb{F}, \forall \mathbf{x} \in L: \lambda\mathbf{x} \in L\) (замкнутость относительно умножения на скаляр).
\end{enumerate}
\end{theorem}

\begin{proof}
Докажем теорему в два этапа.

\textbf{Необходимость.} Если \(L\) -- подпространство, то оно само является линейным пространством. Следовательно, определенные на нем операции сложения и умножения на скаляр не выводят за его пределы, то есть условия 1 и 2 выполняются по определению линейного пространства.

\textbf{Достаточность.} Пусть для непустого множества \(L\) выполняются условия 1 и 2. Покажем, что \(L\) удовлетворяет всем аксиомам линейного пространства.
\begin{itemize}
    \item Надо проверить аксиомы 3 и 4, остальные очевидно следуют из свойств $V$.
    \item Из условия 2, взяв \(\lambda = 0 \in \mathbb{F}\), получаем \(0 \cdot \mathbf{x} = \mathbf{0} \in L\) для любого \(\mathbf{x} \in L\) (используем доказанное ранее свойство \(0\cdot\mathbf{x} = \mathbf{0}\)). Следовательно, нулевой вектор принадлежит \(L\) (аксиома 3).
    \item Из условия 2, взяв \(\lambda = -1 \in \mathbb{F}\), получаем \((-1)\cdot\mathbf{x} = -\mathbf{x} \in L\) (используя свойство \((-1)\mathbf{x} = -\mathbf{x}\)). Таким образом, для каждого вектора из \(L\) противоположный ему также лежит в \(L\) (аксиома 4).
\end{itemize}
\end{proof}

\begin{example}
Рассмотрим пространство \(\mathbb{R}^3\). Множество \(L = \{(x_1, x_2, x_3) \in \mathbb{R}^3 \mid x_1 + x_2 + x_3 = 0\}\) является подпространством. Действительно, если взять два вектора из \(L\), то сумма их координат также будет равна нулю, и при умножении любого вектора на скаляр это свойство сохраняется.
\end{example}

\begin{example}
Множество многочленов степени не выше \(n\) (обозначается \(\mathbb{F}_n[x]\)) является подпространством в пространстве всех многочленов \(\mathbb{F}[x]\).
\end{example}

\section*{Линейная зависимость и независимость системы векторов}

\begin{definition}[Линейная комбинация]
Пусть задана система \(n\) векторов \(\{\mathbf{a}_1, \mathbf{a}_2, \dots, \mathbf{a}_n\}\) и \(n\) действительных чисел \(\lambda_1, \dots, \lambda_n\). Вектор
\begin{equation}
\label{eq:lincomb}
\lambda_1 \mathbf{a}_1 + \lambda_2 \mathbf{a}_2 + \dots + \lambda_n \mathbf{a}_n
\end{equation}
называется \textbf{линейной комбинацией} данной системы векторов.
\end{definition}

\begin{definition}[Линейная независимость]
Система векторов \(\{\mathbf{a}_1, \dots, \mathbf{a}_n\}\) называется \textbf{линейно независимой} (ЛНЗ), если равенство линейной комбинации нулевому вектору:
\[
\lambda_1 \mathbf{a}_1 + \dots + \lambda_n \mathbf{a}_n = \mathbf{0}
\]
возможно только в тривиальном случае, когда все коэффициенты \(\lambda_i = 0\).
\end{definition}

\begin{definition}[Линейная зависимость]
Система векторов \(\{\mathbf{a}_1, \dots, \mathbf{a}_n\}\) называется \textbf{линейно зависимой} (ЛЗ), если существует нетривиальная линейная комбинация, равная нулевому вектору, т.е. существуют коэффициенты \(\lambda_i\), не все равные нулю, для которых выполняется равенство:
\[
\lambda_1 \mathbf{a}_1 + \dots + \lambda_n \mathbf{a}_n = \mathbf{0}.
\]
\end{definition}

\subsection*{Основная теорема о линейной зависимости}

\begin{theorem}\label{th:linind}
Пусть дана система векторов \(\{\mathbf{v}_1, \dots, \mathbf{v}_m\}\). Тогда справедливы следующие утверждения:
\begin{enumerate}
    \item Система линейно зависима тогда и только тогда, когда хотя бы один из векторов этой системы можно представить в виде линейной комбинации остальных.
    \item Если некоторая подсистема данной системы векторов линейно зависима, то и вся система линейно зависима.
    \item\label{pt:addvec} Если векторы \(\mathbf{v}_1, \dots, \mathbf{v}_m\) линейно независимы, а система \(\{\mathbf{v}_1, \dots, \mathbf{v}_m, \mathbf{v}_{m+1}\}\) линейно зависима, то вектор \(\mathbf{v}_{m+1}\) можно представить в виде линейной комбинации предыдущих.
\end{enumerate}
\end{theorem}

\begin{proof}
\textbf{Пункт 1.}
\begin{itemize}
    \item[\(\Rightarrow\)] Пусть система ЛЗ. Тогда существуют числа \(\lambda_1, \dots, \lambda_m\), не все равные нулю, такие что \(\sum_{i=1}^m \lambda_i \mathbf{v}_i = \mathbf{0}\). Пусть \(\lambda_k \neq 0\). Тогда можно выразить \(\mathbf{v}_k\):
    \[
    \lambda_k \mathbf{v}_k = -\sum_{i \neq k} \lambda_i \mathbf{v}_i \quad \Rightarrow \quad \mathbf{v}_k = \sum_{i \neq k} \left(-\frac{\lambda_i}{\lambda_k}\right) \mathbf{v}_i,
    \]
    то есть \(\mathbf{v}_k\) есть линейная комбинация остальных.
    \item[\(\Leftarrow\)] Пусть, например, \(\mathbf{v}_1 = \mu_2 \mathbf{v}_2 + \dots + \mu_m \mathbf{v}_m\). Перенеся все в одну сторону, получим:
    \[
    \mathbf{v}_1 - \mu_2 \mathbf{v}_2 - \dots - \mu_m \mathbf{v}_m = \mathbf{0}.
    \]
    Коэффициент при \(\mathbf{v}_1\) равен 1 (не ноль), следовательно, система линейно зависима.
\end{itemize}

\textbf{Пункт 2.} Пусть подсистема из первых \(k\) векторов (\(k \le m\)) ЛЗ. Тогда существуют \(\lambda_1, \dots, \lambda_k\), не все равные нулю, такие что \(\lambda_1 \mathbf{v}_1 + \dots + \lambda_k \mathbf{v}_k = \mathbf{0}\). Добавив к этой сумме остальные векторы с нулевыми коэффициентами, получим:
\[
\lambda_1 \mathbf{v}_1 + \dots + \lambda_k \mathbf{v}_k + 0\cdot \mathbf{v}_{k+1} + \dots + 0\cdot \mathbf{v}_m = \mathbf{0}.
\]
Так как среди коэффициентов есть ненулевые (из первой группы), то вся система ЛЗ.

\textbf{Пункт 3.} Так как система \(\{\mathbf{v}_1, \dots, \mathbf{v}_{m+1}\}\) ЛЗ, то существуют \(\lambda_1, \dots, \lambda_{m+1}\), не все равные нулю, что \(\sum_{i=1}^{m+1} \lambda_i \mathbf{v}_i = \mathbf{0}\). Если \(\lambda_{m+1} = 0\), то мы получили бы нетривиальную комбинацию первых \(m\) векторов, равную нулю, что противоречит их линейной независимости. Следовательно, \(\lambda_{m+1} \neq 0\). Тогда, аналогично пункту 1, выражаем \(\mathbf{v}_{m+1}\):
\[
\mathbf{v}_{m+1} = -\frac{\lambda_1}{\lambda_{m+1}} \mathbf{v}_1 - \dots - \frac{\lambda_m}{\lambda_{m+1}} \mathbf{v}_m.
\]
\end{proof}

\subsection*{Следствия из теоремы}

\begin{corollary}
\begin{enumerate}
    \item Если система векторов содержит нулевой вектор, то она линейно зависима.
    \begin{proof}
        Пусть \(\mathbf{v}_1 = \mathbf{0}\). Тогда комбинация \(1\cdot \mathbf{v}_1 + 0\cdot \mathbf{v}_2 + \dots + 0\cdot \mathbf{v}_n = \mathbf{0}\) является нетривиальной (коэффициент при \(\mathbf{v}_1\) отличен от нуля). Значит, система ЛЗ.
    \end{proof}

    \item Если система содержит два пропорциональных (или равных) вектора, то она линейно зависима.
    \begin{proof}
        Пусть \(\mathbf{v}_2 = \alpha \mathbf{v}_1\). Тогда \(1 \cdot \mathbf{v}_2 - \alpha \cdot \mathbf{v}_1 = \mathbf{0}\) — нетривиальная комбинация.
    \end{proof}

    \item Если система векторов линейно независима, то любая ее подсистема также линейно независима.
    \begin{proof}
        Это следует из пункта 2 теоремы. Если бы какая-то подсистема была ЛЗ, то и вся система была бы ЛЗ, что противоречит условию.
    \end{proof}
\end{enumerate}
\end{corollary}

\subsection*{Теорема о "прополке" (о максимальной линейно независимой подсистеме)}

\begin{definition}[Линейная оболочка]
\textbf{Линейной оболочкой} системы векторов \(\{\mathbf{a}_1, \dots, \mathbf{a}_n\}\) (обозначается \(\operatorname{span}(\mathbf{a}_1, \dots, \mathbf{a}_n)\)) называется множество всевозможных линейных комбинаций этих векторов:
\[
\operatorname{span}(\mathbf{a}_1, \dots, \mathbf{a}_n) = \left\{ \sum_{i=1}^n \lambda_i \mathbf{a}_i \;\big|\; \lambda_i \in \mathbb{R} \right\}.
\]
Легко проверить, что линейная оболочка сама является линейным подпространством (наименьшим, содержащим данную систему).
\end{definition}

\begin{theorem}
Из любой конечной системы векторов \(\{\mathbf{a}_1, \dots, \mathbf{a}_n\}\), содержащей больше одного вектора, можно выделить линейно независимую подсистему \(\{\mathbf{a}_{i_1}, \dots, \mathbf{a}_{i_k}\}\) такую, что её линейная оболочка совпадает с линейной оболочкой исходной системы:
\[
\operatorname{span}(\mathbf{a}_{i_1}, \dots, \mathbf{a}_{i_k}) = \operatorname{span}(\mathbf{a}_1, \dots, \mathbf{a}_n).
\]
\end{theorem}

\begin{proof}
Будем последовательно просматривать векторы исходной системы и строить искомую подсистему.
\begin{enumerate}
    \item Возьмём первый ненулевой вектор (если все нулевые, то оболочка тривиальна, и можно взять пустую подсистему). Обозначим его \(\mathbf{b}_1\). Система из одного ненулевого вектора линейно независима.
    \item Для каждого следующего вектора \(\mathbf{a}_j\) пытаемся добавить к текущей системе. Если полученная система линейно независима, оставляем его. Если линейно зависима — отбрасываем (выбрасываем), так как по пункту \ref{pt:addvec} теоремы \ref{th:linind} он выражается через предыдущие (т.е. $\mathbf{a}_j \in \operatorname{span}(\mathbf{b}_1, \dots, \mathbf{b}_r)$).
    \item Процесс заканчивается, когда все векторы исходной системы рассмотрены.
\end{enumerate}
Полученная подсистема \(\{\mathbf{b}_1, \dots, \mathbf{b}_k\}\) по построению линейно независима. Каждый вектор исходной системы либо сам попал в подсистему, либо был отброшен как линейная комбинация уже отобранных. Следовательно, все исходные векторы лежат в \(\operatorname{span}(\mathbf{b}_1, \dots, \mathbf{b}_k)\), а значит и их линейная оболочка содержится в этом подпространстве. Обратное включение очевидно, так как каждый \(\mathbf{b}_i\) является одним из исходных векторов. Таким образом, оболочки равны.
\end{proof}

\section*{Базис линейного пространства}

Везде далее \(V\) — линейное пространство над полем \(\mathbb{F}\) (например, \(\mathbb{R}\) или \(\mathbb{C}\)). Рассматриваются только конечные системы векторов.

\begin{definition}[Порождающая система]
Система векторов \(\{\mathbf{a}_1,\dots,\mathbf{a}_n\} \subset V\) называется \textbf{порождающей} (или системой образующих) пространства \(V\), если каждый вектор \(\mathbf{v} \in V\) может быть представлен в виде линейной комбинации векторов этой системы, т.е. \(V = \operatorname{span}(\mathbf{a}_1,\dots,\mathbf{a}_n)\).
\end{definition}

\subsection*{Эквивалентные определения базиса}

\begin{theorem}\label{th:basis}
Для системы векторов \(\{\mathbf{v}_1,\dots,\mathbf{v}_n\}\) в пространстве \(V\) следующие условия эквивалентны:
\begin{enumerate}
    \item \(\{\mathbf{v}_1,\dots,\mathbf{v}_n\}\) — линейно независима и порождает \(V\).
    \item \(\{\mathbf{v}_1,\dots,\mathbf{v}_n\}\) — \textbf{максимальная линейно независимая система}: она линейно независима, и любая линейно независимая система векторов в \(V\) содержит не более \(n\) векторов.
    \item \(\{\mathbf{v}_1,\dots,\mathbf{v}_n\}\) — \textbf{минимальная порождающая система}: она порождает \(V\), и любая порождающая система векторов в \(V\) содержит не менее \(n\) векторов.
\end{enumerate}
\end{theorem}

\begin{proof}
Проведём доказательство по схеме \(1 \Rightarrow 2\), \(2 \Rightarrow 1\), \(1 \Rightarrow 3\), \(3 \Rightarrow 1\).

\paragraph{\(1 \Rightarrow 2\).}
Пусть \(\{\mathbf{v}_1,\dots,\mathbf{v}_n\}\)~--- порождающая. Возьмём произвольную линейно независимую систему \(\{\mathbf{u}_1,\dots,\mathbf{u}_m\}\). Рассмотрим систему векторов \(\{\mathbf{u}_m, \mathbf{v}_1,\dots,\mathbf{v}_n\}\). Она линейно зависима, поскольку \(\{\mathbf{v}_1,\dots,\mathbf{v}_n\}\)~--- порождающая. Применим теорему о <<прополке>>, получим линейно независимую систему \(\{\mathbf{u}_m, \dots,\mathbf{v}_i, \dots\}\), в которой был отброшен хотя бы один вектор $\mathbf{v}$, однако линейная оболочка сохранится и будет равняться $V$, а значит система опять порождающая. На каждом шаге будем добавлять вектор $\mathbf{u}$ к началу порождающей системы из прошлого шага, получать линейно зависимую систему и применять теорему о <<прополке>>. По свойствам алгоритма и линейной независимости системы $\mathbf{u}$, все векторы $\mathbf{u}$ будут оставться в системе, а так же число векторов в системе на каждом шаге не увеличивается, поскольку выкидываем хотя бы один вектор. В результате получим линейно независимую порождающую систему \(\{\mathbf{u}_1, \dots, \mathbf{u}_m, \dots,\mathbf{v}_i, \dots\}\) размера не больше изначального $n$, следовательно $m \leq n$. Заметим, что здесь мы не использовали свойство линейной независимости $\mathbf{v}$.

\paragraph{\(2 \Rightarrow 1\).}
Пусть \(\{\mathbf{v}_1,\dots,\mathbf{v}_n\}\) максимальна по числу элементов среди линейно независимых систем. Докажем, что она порождает \(V\). Возьмём произвольный вектор \(\mathbf{w} \in V\). Рассмотрим систему \(\{\mathbf{v}_1,\dots,\mathbf{v}_n,\mathbf{w}\}\). Если бы \(\mathbf{w}\) не выражался через \(\mathbf{v}_i\), то эта система была бы линейно независимой (иначе из зависимости следовало бы выражение \(\mathbf{w}\) через \(\mathbf{v}_i\)). Но тогда мы получили бы линейно независимую систему из \(n+1\) вектора, что противоречит максимальности. Значит, \(\mathbf{w}\) выражается через \(\mathbf{v}_i\). Таким образом, система порождает \(V\).

\paragraph{\(1 \Rightarrow 3\).}
\(\{\mathbf{v}_1,\dots,\mathbf{v}_n\}\)~--- линейно независимая. Пусть \(\{\mathbf{u}_1,\dots,\mathbf{u}_m\}\)~--- порождающая. Применим пункт \(1 \Rightarrow 2\) для систем $\mathbf{u}$ как порождающей и $\mathbf{v}$ как линейно независимой (поменяем местами $\mathbf{v}$ и $\mathbf{u}$). Получим, что $n \leq m$.

\paragraph{\(3 \Rightarrow 1\).}
Пусть \(\{\mathbf{v}_1,\dots,\mathbf{v}_n\}\) минимальна по числу элементов среди порождающих систем. Докажем, что она линейно независима. Предположим противное: система линейно зависима. Тогда один из векторов, скажем \(\mathbf{v}_1\), выражается через остальные:
\[
\mathbf{v}_1 = \gamma_2\mathbf{v}_2 + \dots + \gamma_n\mathbf{v}_n.
\]
Удалим \(\mathbf{v}_1\). Оставшаяся система \(\{\mathbf{v}_2,\dots,\mathbf{v}_n\}\) также порождает \(V\), что противоречит минимальности. Следовательно, система линейно независима.
\end{proof}

\begin{definition}[Базис]
\textbf{Базисом} линейного пространства \(V\) называется любая система векторов, которая удовлетворяет любому из условий теоремы \ref{th:basis}.

Размерностью линейного пространства $\dim V$ называется число элементов базиса $V$ (что равняется максимальному размеру линейно независимой системы и минимальному размеру порождающей системы векторов).
\end{definition}

\subsection*{Два важных свойства базиса в конечномерном пространстве}

Предположим, что пространство \(V\) конечномерно (т.е. имеет конечный базис). Тогда справедливы следующие утверждения.

\begin{theorem}[О дополнении до базиса]
Любую линейно независимую систему векторов \(\{\mathbf{u}_1,\dots,\mathbf{u}_k\}\) в конечномерном пространстве \(V\) можно дополнить до базиса \(V\).
\end{theorem}

\begin{proof}
Если \(\operatorname{span}(\mathbf{u}_1,\dots,\mathbf{u}_k) = V\), то система уже является базисом. Иначе существует вектор \(\mathbf{u}_{k+1} \notin \operatorname{span}(\mathbf{u}_1,\dots,\mathbf{u}_k)\). Добавим его — полученная система \(\{\mathbf{u}_1,\dots,\mathbf{u}_k,\mathbf{u}_{k+1}\}\) остаётся линейно независимой. Продолжаем этот процесс. Так как пространство конечномерно, и максимальный размер линейно независимой системы векторов ограничен размерностью пространства, процесс остановится через конечное число шагов.
\end{proof}

\begin{theorem}[О выделении базиса из порождающей системы]
Из любой конечной порождающей системы векторов пространства \(V\) можно выделить подсистему, являющуюся базисом \(V\).
\end{theorem}

\begin{proof}
Следует из теоремы о <<прополке>>.
\end{proof}
\section{Координаты вектора в заданном базисе}

Пусть $V$ — линейное пространство над полем $\R$, и пусть $e = (e_1, e_2, \dots, e_n)$ — базис $V$. Тогда любой вектор $x \in V$ единственным образом представляется в виде линейной комбинации
\[
x = x^1 e_1 + x^2 e_2 + \dots + x^n e_n.
\]

\begin{proof}[Доказательство единственности]
Предположим, что существуют два представления:
\[
x = \sum_{i=1}^n x^i e_i = \sum_{i=1}^n y^i e_i.
\]
Тогда $\sum_{i=1}^n (x^i - y^i) e_i = 0$. Так как система $e_1,\dots,e_n$ линейно независима, все коэффициенты $x^i - y^i = 0$, т.е. $x^i = y^i$.
\end{proof}

Числа $x^i$ называются \textbf{координатами} вектора $x$ в базисе $e$. Координатный столбец обозначается $[x]_e$.

\section{Переход к новому базису}

Пусть $e = (e_1,\dots,e_n)$ и $f = (f_1,\dots,f_n)$ — два базиса $V$. Тогда каждый вектор $f_j$ выражается через базис $e$:
\[
f_j = \sum_{i=1}^n C_{ij}\, e_i, \quad j=1,\dots,n.
\]
Матрица $C = (C_{ij})_{n\times n}$ называется \textbf{матрицей перехода} от базиса $e$ к базису $f$.

\begin{lemma}
Матрица перехода $C$ невырождена.
\end{lemma}
\begin{proof}
Если бы $\det C = 0$, то столбцы $C$ были бы линейно зависимы, т.е. существовал бы ненулевой набор $\alpha_1,\dots,\alpha_n$ такой, что $\sum_{j} C_{ij}\alpha_j = 0$ для всех $i$. Тогда $\sum_j \alpha_j f_j = \sum_j \alpha_j \sum_i C_{ij} e_i = \sum_i \bigl(\sum_j C_{ij}\alpha_j\bigr) e_i = 0$, что противоречит линейной независимости $f_j$. Следовательно, $C$ обратима.
\end{proof}

\textbf{Преобразование координат.} Пусть $x \in V$ имеет координаты $[x]_e = (x^1,\dots,x^n)^T$ в базисе $e$ и $[x]_f = (x'^1,\dots,x'^n)^T$ в базисе $f$. Тогда
\[
[x]_e = C\, [x]_f,\qquad [x]_f = C^{-1} [x]_e.
\]

\begin{proof}
Имеем $x = \sum_{j} x'^j f_j = \sum_{j} x'^j \sum_i C_{ij} e_i = \sum_i \bigl( \sum_j C_{ij} x'^j \bigr) e_i$. Следовательно, координата при $e_i$ равна $\sum_j C_{ij} x'^j$, что и означает $[x]_e = C [x]_f$.
\end{proof}

\section{Скалярное произведение}

\textbf{Определение.} Скалярным произведением в вещественном линейном пространстве $V$ называется функция $(\cdot,\cdot): V\times V \to \R$, удовлетворяющая:
\begin{enumerate}
  \item линейность по первому аргументу: $(\lambda x + y, z) = \lambda (x,z) + (y,z)$;
  \item симметричность: $(x,y) = (y,x)$;
  \item положительная определённость: $(x,x) \ge 0$, причём $(x,x)=0 \iff x=0$.
\end{enumerate}
Пространство с таким произведением называется \textbf{евклидовым}.

Важнейшее следствие — неравенство Коши--Буняковского:
\[
|(x,y)| \le \sqrt{(x,x)}\sqrt{(y,y)}.
\]

\begin{proof}
Для любого $t \in \R$ рассмотрим $(x+ty, x+ty) \ge 0$. Раскрывая: $(x,x) + 2t(x,y) + t^2(y,y) \ge 0$. Это квадратный трёхчлен относительно $t$, неотрицательный при всех $t$, значит его дискриминант $\le 0$:
\[
(2(x,y))^2 - 4(x,x)(y,y) \le 0 \quad\Rightarrow\quad (x,y)^2 \le (x,x)(y,y).
\]
Извлекая корень, получаем требуемое.
\end{proof}

\subsection{Норма вектора}

\begin{definition}
    Пусть $V$ — линейное пространство над $\mathbb{R}$ (или $\mathbb{C}$). \textbf{Нормой} на $V$ называется функция $\|\cdot\|: V \to \mathbb{R}$, удовлетворяющая для любых $x,y\in V$ и любого скаляра $\lambda$ следующим аксиомам:
    \begin{enumerate}
        \item \textbf{Положительная определённость:} $\|x\| \ge 0$, причём $\|x\| = 0$ тогда и только тогда, когда $x = \mathbf{0}$.
        \item \textbf{Абсолютная однородность:} $\|\lambda x\| = |\lambda| \cdot \|x\|$.
        \item \textbf{Неравенство треугольника:} $\|x + y\| \le \|x\| + \|y\|$.
    \end{enumerate}
    Пространство с заданной на нём нормой называется \textbf{нормированным пространством}.
\end{definition}

\begin{theorem}[Норма, индуцированная скалярным произведением]
    В евклидовом пространстве $V$ функция $\|x\| = \sqrt{(x,x)}$ удовлетворяет всем аксиомам нормы.
\end{theorem}
\begin{proof}
    \begin{enumerate}
        \item Положительная определённость следует из положительной определённости скалярного произведения: $(x,x) \ge 0$, и $(x,x)=0 \Leftrightarrow x=\mathbf{0}$.
        \item Абсолютная однородность: $\|\lambda x\| = \sqrt{(\lambda x,\lambda x)} = \sqrt{\lambda\overline{\lambda}\,(x,x)} = |\lambda|\,\|x\|$ (в вещественном случае $\overline{\lambda}=\lambda$).
        \item Неравенство треугольника является следствием неравенства Коши–Буняковского:
        \[
        \|x+y\|^2 = (x+y,x+y) = \|x\|^2 + \|y\|^2 + 2(x,y) \le \|x\|^2 + \|y\|^2 + 2\|x\|\,\|y\| = (\|x\|+\|y\|)^2,
        \]
        откуда $\|x+y\| \le \|x\|+\|y\|$ (извлекая квадратный корень из обеих частей).
    \end{enumerate}
\end{proof}

\section{Ортогональные и ортонормированные системы}

\textbf{Определение.} Векторы $x$ и $y$ ортогональны, если $(x,y)=0$. Система векторов $\{v_1,\dots,v_k\}$ называется \textbf{ортогональной}, если $(v_i, v_j)=0$ при $i\neq j$. Если к тому же $\|v_i\|=1$ для всех $i$, система \textbf{ортонормирована}.

\begin{theorem}
Любая ортогональная система ненулевых векторов линейно независима.
\end{theorem}
\begin{proof}
Пусть $\sum_{i=1}^k \alpha_i v_i = 0$. Умножим скалярно на $v_j$ (для фиксированного $j$):
\[
0 = \Bigl(\sum_i \alpha_i v_i, v_j\Bigr) = \sum_i \alpha_i (v_i, v_j) = \alpha_j (v_j, v_j),
\]
так как $(v_i, v_j)=0$ при $i\neq j$. Поскольку $v_j \neq 0$, имеем $(v_j, v_j)>0$, следовательно $\alpha_j = 0$. Это верно для всех $j$, значит все коэффициенты нулевые.
\end{proof}

\section{Процесс ортогонализации Грама–Шмидта}

Пусть $\{a_1,\dots,a_n\}$ — линейно независимая система в евклидовом пространстве. Построим по ней ортогональную систему $\{b_1,\dots,b_n\}$ и ортонормированную $\{u_1,\dots,u_n\}$.

\begin{enumerate}
  \item Положим $b_1 = a_1$.
  \item Для $k = 2,\dots,n$ последовательно вычисляем
  \[
  b_k = a_k - \sum_{j=1}^{k-1} \frac{(a_k, b_j)}{(b_j, b_j)}\, b_j.
  \]
\end{enumerate}

\begin{proof}[Доказательство ортогональности (индукция)]
База: для $k=1$ тривиально. Предположим, что $b_1,\dots,b_{k-1}$ попарно ортогональны. Проверим, что $b_k$ ортогонален каждому $b_i$ при $i<k$.
\[
(b_k, b_i) = (a_k, b_i) - \sum_{j=1}^{k-1} \frac{(a_k, b_j)}{(b_j, b_j)} (b_j, b_i).
\]
В силу индукционного предположения $(b_j, b_i)=0$ при $j\neq i$, а при $j=i$ получаем $(b_i, b_i)$. Тогда сумма сводится к одному слагаемому:
\[
(b_k, b_i) = (a_k, b_i) - \frac{(a_k, b_i)}{(b_i, b_i)}(b_i, b_i) = 0.
\]
Таким образом, все $b_1,\dots,b_k$ ортогональны.
\end{proof}

После построения ортогональной системы $\{b_i\}$ нормируем её: $u_i = \frac{b_i}{\|b_i\|}$. Полученная система $\{u_i\}$ ортонормирована и, очевидно, является базисом (поскольку количество векторов равно размерности и они линейно независимы). Тем самым доказано существование ортонормированного базиса в любом конечномерном евклидовом пространстве.

\section{Скалярное произведение в координатах}

Пусть $V$ — евклидово пространство со скалярным произведением $(\cdot,\cdot)$, и пусть $e = (e_1,\dots,e_n)$ — произвольный базис $V$. Для любых векторов
\[
x = \sum_{i=1}^n x^i e_i, \qquad y = \sum_{j=1}^n y^j e_j
\]
их скалярное произведение выражается через координаты следующим образом:
\[
(x,y) = \sum_{i=1}^n\sum_{j=1}^n x^i y^j (e_i, e_j).
\]

\begin{definition}
    \textbf{Матрицей Грама} базиса $e$ называется матрица
    \[
    G(e) = \bigl( (e_i, e_j) \bigr)_{i,j=1}^n.
    \]
    Тогда для координатных столбцов $[x]_e = (x^1,\dots,x^n)^T$ и $[y]_e = (y^1,\dots,y^n)^T$ имеем
    \[
    (x,y) = [x]_e^T \, G(e) \, [y]_e.
    \]
\end{definition}

\begin{remark}
    Матрица Грама всегда симметрична ($G^T = G$) и положительно определена (поскольку $(x,x) > 0$ для любого ненулевого $x$). Её определитель называется определителем Грама и равен нулю тогда и только тогда, когда векторы линейно зависимы.
\end{remark}

\subsection{Скалярное произведение в ортонормированном базисе}

В ортонормированном базисе матрица Грама становится единичной: $G(e) = I_n$. Поэтому скалярное произведение принимает простейший вид:
\[
(x,y) = \sum_{i=1}^n x^i y^i = [x]_e^T [y]_e.
\]

Кроме того, в ортонормированном базисе координаты вектора совпадают с его скалярными произведениями на базисные векторы:
\[
x^i = (x, e_i) \quad \text{для всех } i = 1,\dots,n.
\]
\begin{proof}
    Подставим $x = \sum_{j} x^j e_j$ в $(x, e_i)$:
    \[
    (x, e_i) = \sum_{j} x^j (e_j, e_i) = \sum_{j} x^j \delta_{ji} = x^i.
    \]
\end{proof}

\subsection{Преобразование скалярного произведения при замене базиса}

Пусть $e$ и $f$ — два базиса, связанные матрицей перехода $T$: $f = e T$ (т.е. $f_j = \sum_i T_{ij} e_i$). Тогда для координатных столбцов вектора $x$ в этих базисах справедливо $[x]_e = T [x]_f$. Выразим скалярное произведение через координаты в базисе $f$:
\[
(x,y) = [x]_e^T G(e) [y]_e = (T[x]_f)^T G(e) (T[y]_f) = [x]_f^T \bigl( T^T G(e) T \bigr) [y]_f.
\]
Следовательно, матрица Грама в базисе $f$ равна
\[
G(f) = T^T G(e) T.
\]
Если оба базиса ортонормированны, то $G(e)=I$, $G(f)=I$, и $T^T T = I$.

\end{document}
